% Generated draft: requires mathematical and source review.
\documentclass[11pt]{article}
\usepackage[margin=25mm]{geometry}
\usepackage{fontspec}
\setmainfont{Latin Modern Roman}
\usepackage{amsmath,amssymb,mathtools,unicode-math}
\setmathfont{Latin Modern Math}
\usepackage{enumitem,needspace,xurl,hyperref}
\setlength{\parindent}{0pt}
\setlength{\parskip}{5pt}
\setlength{\emergencystretch}{3em}
\newcommand{\N}{\mathbb{N}}
\newcommand{\Z}{\mathbb{Z}}
\newcommand{\Q}{\mathbb{Q}}
\newcommand{\R}{\mathbb{R}}
\newcommand{\C}{\mathbb{C}}
\newcommand{\F}{\mathbb{F}}
\newcommand{\A}{\mathbb{A}}
\newcommand{\PP}{\mathbb{P}}
\newcommand{\E}{\mathbb{E}}
\newcommand{\eps}{\varepsilon}
\newcommand{\dd}{\,\mathrm{d}}
\begin{document}
\section*{UnsolvedMath ID: 20000730}
\textbf{Source problem code:} catalogue entry
\subsection{Definitions and mathematical statement}
The source asks for relations between modularity of a series and its expansion in \(q\)-orthogonal polynomials. It specifies neither a polynomial family, expansion coefficients, modular group nor weight, so it does not assert that every such expansion is modular.

The supplied special case uses \(q=e^{2\pi i\tau}\), \(\operatorname{Im}\tau>0\), with \(q^{1/2}=e^{\pi i\tau}\), and continuous \(q\)-Hermite polynomials normalized by \(H_0=1\), \(H_1=2x\), and
\[
 H_{n+1}(x\mid q)=2xH_n(x\mid q)-(1-q^n)H_{n-1}(x\mid q).
\]
Write \((a;q)_n=\prod_{j=0}^{n-1}(1-aq^j)\) and use the corresponding infinite product. For \(\beta\in\Q\), \(\zeta=e^{2\pi i\beta}\), define
\[
 F_\beta(\tau)=q^{1/24}\sum_{n\ge0}
 \frac{H_n(\cos(2\pi\beta)\mid q)q^{n/2}}{(q;q)_n}.
\]
The report asserts
\[
 F_\beta(\tau)=\frac{q^{1/24}}
 {(\zeta q^{1/2};q)_\infty(\zeta^{-1}q^{1/2};q)_\infty}
 =-e^{-\pi i\beta/2}\,g_{(1/2,\beta)}(\tau)^{-1}.
\]
Its Siegel-function normalization is fixed by
\[
 g_{(1/2,\beta)}=-e^{-\pi i\beta/2}q^{-1/24}
 (\zeta q^{1/2};q)_\infty(\zeta^{-1}q^{1/2};q)_\infty.
\]
If \(M\) is the least positive integer with \(M/2,M\beta\in\Z\), the report gives \(F_\beta^{12M}\) as a modular unit on \(X(M)\). A modular unit is a meromorphic modular function with zeros and poles confined to cusps; this power statement is distinct from claiming \(F_\beta\) has trivial multiplier itself.

The reported logarithmic derivative is a weight-two modular form for the corresponding level convention and has expansion
\[
 \frac1{2\pi i}\frac{d}{d\tau}\log F_\beta
 =\frac1{24}+\sum_{m\ge1}\left[\frac12
 \sum_{\substack{d\mid m\\d\text{ odd}}}
 d\bigl(\zeta^{m/d}+\zeta^{-m/d}\bigr)\right]q^{m/2}.
\]
The scalar multiplier and logarithm branch do not affect the derivative. The report emphasizes that mere existence of an orthogonal-polynomial expansion cannot certify modularity without extra arithmetic structure; here rational torsion specialization and the explicit Siegel-product identity provide that structure. It supplies no universal criterion for arbitrary \(q\)-orthogonal families or coefficients.
\subsection{Short English statement}
Seek connections between \(q\)-orthogonal expansions and modularity. The supplied report gives a rationally specialized continuous \(q\)-Hermite series equal to an inverse Siegel function, with an explicit modular-unit power and derivative coefficients. It does not turn arbitrary orthogonal expansions into modularity criteria.
\subsection{Sources}
\begin{itemize}
\item Local UnsolvedMath registry, permanent ID \(20000730\), source code \texttt{catalogue entry}; AIM analytic number theory corpus, \texttt{aim-analytic-number-theory-notes.json}, notes[93]. Attribution: Ulam AI / UnsolvedMath Contributors, CC BY 4.0.
\item AIM \emph{Mock modular forms in combinatorics and arithmetic geometry}, March 8--12, 2010, Problem 1.3: \url{https://aimath.org/WWN/mockmodular/mockmodular.pdf}. Original remarks and literature fields are empty.
\item Registry dated August 14, 2026: statement exact, status \texttt{partially\_solved}, classification \texttt{PARTIAL-PROGRESS}; attempt 1 labelled partial result, rigor proved, novelty confidence low. The torsion \(q\)-Hermite/Siegel identity and derivative certificate are machine-generated research aids, not peer reviewed, requiring independent verification before citation. No full primary references or proofs are supplied.
\end{itemize}
\end{document}
